% paper32-sequence-encodings.tex — Paper 32 of the Judgment Geometry series
% "Sequence and Mutation Encodings: Verifying Python Lists, Dicts, and Mutable State"
\documentclass[11pt]{article}
% jugeo-common.tex — shared preamble for all Judgment Geometry papers
% Include via: % jugeo-common.tex — shared preamble for all Judgment Geometry papers
% Include via: % jugeo-common.tex — shared preamble for all Judgment Geometry papers
% Include via: \input{jugeo-common}

% ── Packages ────────────────────────────────────────────────────────────
\usepackage{amsmath,amssymb,amsthm,mathtools}
\usepackage{tikz-cd}
\usepackage{listings}
\usepackage{xcolor}
\usepackage{xspace}
\usepackage{booktabs}
\usepackage{multirow}
\usepackage{enumitem}
\usepackage[expansion=false]{microtype}
\usepackage{hyperref}
\usepackage{cleveref}
% \usepackage{thmtools}  % disabled: not available in this TeX installation
\usepackage{float}
\usepackage{caption}
\usepackage{subcaption}
\usepackage{tabularx}
\usepackage{stmaryrd}       % \llbracket, \rrbracket
\usepackage{bbm}            % \mathbbm{1}

% ── Page geometry ───────────────────────────────────────────────────────
\usepackage[margin=1in]{geometry}

% ── Theorem environments ────────────────────────────────────────────────
\theoremstyle{plain}
\newtheorem{theorem}{Theorem}[section]
\newtheorem{lemma}[theorem]{Lemma}
\newtheorem{proposition}[theorem]{Proposition}
\newtheorem{corollary}[theorem]{Corollary}

\theoremstyle{definition}
\newtheorem{definition}[theorem]{Definition}
\newtheorem{example}[theorem]{Example}
\newtheorem{construction}[theorem]{Construction}

\theoremstyle{remark}
\newtheorem{remark}[theorem]{Remark}
\newtheorem{notation}[theorem]{Notation}

% ── Python listings style ──────────────────────────────────────────────
\definecolor{jugeoBlue}{HTML}{2563EB}
\definecolor{jugeoGreen}{HTML}{059669}
\definecolor{jugeoRed}{HTML}{DC2626}
\definecolor{jugeoPurple}{HTML}{7C3AED}
\definecolor{jugeoGray}{HTML}{6B7280}
\definecolor{jugeoBg}{HTML}{F8FAFC}

\lstdefinestyle{jugeo-python}{
  language=Python,
  basicstyle=\ttfamily\small,
  keywordstyle=\color{jugeoBlue}\bfseries,
  stringstyle=\color{jugeoGreen},
  commentstyle=\color{jugeoGray}\itshape,
  numberstyle=\tiny\color{jugeoGray},
  backgroundcolor=\color{jugeoBg},
  numbers=left,
  numbersep=6pt,
  frame=single,
  framerule=0.4pt,
  rulecolor=\color{jugeoGray!40},
  breaklines=true,
  breakatwhitespace=true,
  showstringspaces=false,
  tabsize=4,
  xleftmargin=1.5em,
  framexleftmargin=1.5em,
  morekeywords={dataclass,frozen,field,Enum,IntEnum,auto,Any,
                Mapping,Sequence,Optional,match,case},
  literate={→}{{$\to$}}1 {←}{{$\leftarrow$}}1
           {≤}{{$\leq$}}1 {≥}{{$\geq$}}1 {≠}{{$\neq$}}1
           {∈}{{$\in$}}1 {∀}{{$\forall$}}1 {∃}{{$\exists$}}1
           {⊕}{{$\oplus$}}1 {⊖}{{$\ominus$}}1,
  escapeinside={(*@}{@*)},
}
\lstset{style=jugeo-python}

\lstdefinestyle{jugeo-output}{
  basicstyle=\ttfamily\small,
  backgroundcolor=\color{jugeoBg},
  frame=single,
  framerule=0.4pt,
  rulecolor=\color{jugeoGray!40},
  breaklines=true,
  showstringspaces=false,
  xleftmargin=1.5em,
  framexleftmargin=0.5em,
  numbers=none,
}

% ── JuGeo-specific macros ──────────────────────────────────────────────

% Judgment 8-tuple
\newcommand{\J}{\mathcal{J}}
\newcommand{\Jud}[8]{(#1,\,#2,\,#3,\,#4,\,#5,\,#6,\,#7,\,#8)}
\newcommand{\JudShort}[1]{\J_{#1}}

% Components
\newcommand{\coord}{\mathsf{c}}            % coordinate
\newcommand{\prop}{\varphi}                 % proposition
\newcommand{\carrier}{\mathcal{A}}          % carrier type
\newcommand{\evid}{\mathcal{E}}             % evidence bundle
\newcommand{\oblig}{\mathcal{O}}            % obligations
\newcommand{\obstr}{\mathcal{B}}            % obstructions
\newcommand{\trust}{\mathcal{T}}            % trust annotation
\newcommand{\prov}{\Pi}                     % provenance

% Trust algebra
\newcommand{\Talg}{\mathfrak{T}}            % trust ordered algebra
\newcommand{\Eadm}{\mathcal{E}_{\mathrm{adm}}}
\newcommand{\tleq}{\preceq}                 % trust ordering
\newcommand{\tjoin}{\oplus}                 % conservative join
\newcommand{\tatten}{\ominus}               % attenuation
\newcommand{\tpromote}{\uparrow_\pi}        % promotion
\newcommand{\tdemote}{\downarrow_\chi}       % demotion

% Trust tiers
\newcommand{\Tcontra}{\ensuremath{\mathsf{CONTRADICTED}}}
\newcommand{\Tunver}{\ensuremath{\mathsf{UNVERIFIED}}}
\newcommand{\Tcopilot}{\ensuremath{\mathsf{COPILOT}}}
\newcommand{\Toracle}{\ensuremath{\mathsf{ORACLE}}}
\newcommand{\Truntime}{\ensuremath{\mathsf{RUNTIME}}}
\newcommand{\Tsolver}{\ensuremath{\mathsf{SOLVER}}}
\newcommand{\Tproof}{\ensuremath{\mathsf{PROOF}}}

% Site and geometry
\newcommand{\Site}{\mathcal{S}}             % semantic site
\newcommand{\Cov}{\mathrm{Cov}}            % covering family
\newcommand{\Ob}{\mathrm{Ob}}              % objects
\newcommand{\Mor}{\mathrm{Mor}}            % morphisms
\newcommand{\res}{\mathrm{res}}            % restriction
\newcommand{\Cech}{\check{C}}              % Čech complex
\newcommand{\Hcoh}[1]{H^{#1}}             % cohomology group

% Descent
\newcommand{\descent}{\mathrm{desc}}
\newcommand{\glue}{\mathrm{glue}}
\newcommand{\overlap}[2]{#1 \cap #2}

% Semantic moves
\newcommand{\VERIFY}{\textsc{Verify}}
\newcommand{\CONSTRUCT}{\textsc{Construct}}
\newcommand{\NORMALIZE}{\textsc{Normalize}}
\newcommand{\GLUE}{\textsc{Glue}}
\newcommand{\RESTRICT}{\textsc{Restrict}}
\newcommand{\TRANSPORT}{\textsc{Transport}}
\newcommand{\REPAIR}{\textsc{Repair}}
\newcommand{\PROMOTE}{\textsc{Promote}}
\newcommand{\CHALLENGE}{\textsc{Challenge}}
\newcommand{\ESCALATE}{\textsc{Escalate}}

% Fragments
\newcommand{\QFLIA}{\texttt{QF\_LIA}}
\newcommand{\QFLRA}{\texttt{QF\_LRA}}
\newcommand{\QFBV}{\texttt{QF\_BV}}
\newcommand{\QFUF}{\texttt{QF\_UF}}

% Misc
\newcommand{\Python}{\textsc{Python}}
\newcommand{\JuGeo}{\textsc{JuGeo}}
\newcommand{\LEAN}{\textsc{Lean}\,4}
\newcommand{\Fstar}{F$^{\star}$}
\newcommand{\Coq}{\textsc{Coq}}
\newcommand{\Dafny}{\textsc{Dafny}}

% Common shortcuts
\newcommand{\ie}{i.e.\@\xspace}
\newcommand{\eg}{e.g.\@\xspace}
\newcommand{\cf}{cf.\@\xspace}
\newcommand{\wrt}{w.r.t.\@\xspace}

% Evaluation shortcuts
\newcommand{\numcases}{300}
\newcommand{\numspec}{100}
\newcommand{\numeq}{100}
\newcommand{\numbug}{100}
\newcommand{\medtime}{6.5\,ms}
\newcommand{\totaltime}{1.949\,s}


% ── Packages ────────────────────────────────────────────────────────────
\usepackage{amsmath,amssymb,amsthm,mathtools}
\usepackage{tikz-cd}
\usepackage{listings}
\usepackage{xcolor}
\usepackage{xspace}
\usepackage{booktabs}
\usepackage{multirow}
\usepackage{enumitem}
\usepackage[expansion=false]{microtype}
\usepackage{hyperref}
\usepackage{cleveref}
% \usepackage{thmtools}  % disabled: not available in this TeX installation
\usepackage{float}
\usepackage{caption}
\usepackage{subcaption}
\usepackage{tabularx}
\usepackage{stmaryrd}       % \llbracket, \rrbracket
\usepackage{bbm}            % \mathbbm{1}

% ── Page geometry ───────────────────────────────────────────────────────
\usepackage[margin=1in]{geometry}

% ── Theorem environments ────────────────────────────────────────────────
\theoremstyle{plain}
\newtheorem{theorem}{Theorem}[section]
\newtheorem{lemma}[theorem]{Lemma}
\newtheorem{proposition}[theorem]{Proposition}
\newtheorem{corollary}[theorem]{Corollary}

\theoremstyle{definition}
\newtheorem{definition}[theorem]{Definition}
\newtheorem{example}[theorem]{Example}
\newtheorem{construction}[theorem]{Construction}

\theoremstyle{remark}
\newtheorem{remark}[theorem]{Remark}
\newtheorem{notation}[theorem]{Notation}

% ── Python listings style ──────────────────────────────────────────────
\definecolor{jugeoBlue}{HTML}{2563EB}
\definecolor{jugeoGreen}{HTML}{059669}
\definecolor{jugeoRed}{HTML}{DC2626}
\definecolor{jugeoPurple}{HTML}{7C3AED}
\definecolor{jugeoGray}{HTML}{6B7280}
\definecolor{jugeoBg}{HTML}{F8FAFC}

\lstdefinestyle{jugeo-python}{
  language=Python,
  basicstyle=\ttfamily\small,
  keywordstyle=\color{jugeoBlue}\bfseries,
  stringstyle=\color{jugeoGreen},
  commentstyle=\color{jugeoGray}\itshape,
  numberstyle=\tiny\color{jugeoGray},
  backgroundcolor=\color{jugeoBg},
  numbers=left,
  numbersep=6pt,
  frame=single,
  framerule=0.4pt,
  rulecolor=\color{jugeoGray!40},
  breaklines=true,
  breakatwhitespace=true,
  showstringspaces=false,
  tabsize=4,
  xleftmargin=1.5em,
  framexleftmargin=1.5em,
  morekeywords={dataclass,frozen,field,Enum,IntEnum,auto,Any,
                Mapping,Sequence,Optional,match,case},
  literate={→}{{$\to$}}1 {←}{{$\leftarrow$}}1
           {≤}{{$\leq$}}1 {≥}{{$\geq$}}1 {≠}{{$\neq$}}1
           {∈}{{$\in$}}1 {∀}{{$\forall$}}1 {∃}{{$\exists$}}1
           {⊕}{{$\oplus$}}1 {⊖}{{$\ominus$}}1,
  escapeinside={(*@}{@*)},
}
\lstset{style=jugeo-python}

\lstdefinestyle{jugeo-output}{
  basicstyle=\ttfamily\small,
  backgroundcolor=\color{jugeoBg},
  frame=single,
  framerule=0.4pt,
  rulecolor=\color{jugeoGray!40},
  breaklines=true,
  showstringspaces=false,
  xleftmargin=1.5em,
  framexleftmargin=0.5em,
  numbers=none,
}

% ── JuGeo-specific macros ──────────────────────────────────────────────

% Judgment 8-tuple
\newcommand{\J}{\mathcal{J}}
\newcommand{\Jud}[8]{(#1,\,#2,\,#3,\,#4,\,#5,\,#6,\,#7,\,#8)}
\newcommand{\JudShort}[1]{\J_{#1}}

% Components
\newcommand{\coord}{\mathsf{c}}            % coordinate
\newcommand{\prop}{\varphi}                 % proposition
\newcommand{\carrier}{\mathcal{A}}          % carrier type
\newcommand{\evid}{\mathcal{E}}             % evidence bundle
\newcommand{\oblig}{\mathcal{O}}            % obligations
\newcommand{\obstr}{\mathcal{B}}            % obstructions
\newcommand{\trust}{\mathcal{T}}            % trust annotation
\newcommand{\prov}{\Pi}                     % provenance

% Trust algebra
\newcommand{\Talg}{\mathfrak{T}}            % trust ordered algebra
\newcommand{\Eadm}{\mathcal{E}_{\mathrm{adm}}}
\newcommand{\tleq}{\preceq}                 % trust ordering
\newcommand{\tjoin}{\oplus}                 % conservative join
\newcommand{\tatten}{\ominus}               % attenuation
\newcommand{\tpromote}{\uparrow_\pi}        % promotion
\newcommand{\tdemote}{\downarrow_\chi}       % demotion

% Trust tiers
\newcommand{\Tcontra}{\ensuremath{\mathsf{CONTRADICTED}}}
\newcommand{\Tunver}{\ensuremath{\mathsf{UNVERIFIED}}}
\newcommand{\Tcopilot}{\ensuremath{\mathsf{COPILOT}}}
\newcommand{\Toracle}{\ensuremath{\mathsf{ORACLE}}}
\newcommand{\Truntime}{\ensuremath{\mathsf{RUNTIME}}}
\newcommand{\Tsolver}{\ensuremath{\mathsf{SOLVER}}}
\newcommand{\Tproof}{\ensuremath{\mathsf{PROOF}}}

% Site and geometry
\newcommand{\Site}{\mathcal{S}}             % semantic site
\newcommand{\Cov}{\mathrm{Cov}}            % covering family
\newcommand{\Ob}{\mathrm{Ob}}              % objects
\newcommand{\Mor}{\mathrm{Mor}}            % morphisms
\newcommand{\res}{\mathrm{res}}            % restriction
\newcommand{\Cech}{\check{C}}              % Čech complex
\newcommand{\Hcoh}[1]{H^{#1}}             % cohomology group

% Descent
\newcommand{\descent}{\mathrm{desc}}
\newcommand{\glue}{\mathrm{glue}}
\newcommand{\overlap}[2]{#1 \cap #2}

% Semantic moves
\newcommand{\VERIFY}{\textsc{Verify}}
\newcommand{\CONSTRUCT}{\textsc{Construct}}
\newcommand{\NORMALIZE}{\textsc{Normalize}}
\newcommand{\GLUE}{\textsc{Glue}}
\newcommand{\RESTRICT}{\textsc{Restrict}}
\newcommand{\TRANSPORT}{\textsc{Transport}}
\newcommand{\REPAIR}{\textsc{Repair}}
\newcommand{\PROMOTE}{\textsc{Promote}}
\newcommand{\CHALLENGE}{\textsc{Challenge}}
\newcommand{\ESCALATE}{\textsc{Escalate}}

% Fragments
\newcommand{\QFLIA}{\texttt{QF\_LIA}}
\newcommand{\QFLRA}{\texttt{QF\_LRA}}
\newcommand{\QFBV}{\texttt{QF\_BV}}
\newcommand{\QFUF}{\texttt{QF\_UF}}

% Misc
\newcommand{\Python}{\textsc{Python}}
\newcommand{\JuGeo}{\textsc{JuGeo}}
\newcommand{\LEAN}{\textsc{Lean}\,4}
\newcommand{\Fstar}{F$^{\star}$}
\newcommand{\Coq}{\textsc{Coq}}
\newcommand{\Dafny}{\textsc{Dafny}}

% Common shortcuts
\newcommand{\ie}{i.e.\@\xspace}
\newcommand{\eg}{e.g.\@\xspace}
\newcommand{\cf}{cf.\@\xspace}
\newcommand{\wrt}{w.r.t.\@\xspace}

% Evaluation shortcuts
\newcommand{\numcases}{300}
\newcommand{\numspec}{100}
\newcommand{\numeq}{100}
\newcommand{\numbug}{100}
\newcommand{\medtime}{6.5\,ms}
\newcommand{\totaltime}{1.949\,s}


% ── Packages ────────────────────────────────────────────────────────────
\usepackage{amsmath,amssymb,amsthm,mathtools}
\usepackage{tikz-cd}
\usepackage{listings}
\usepackage{xcolor}
\usepackage{xspace}
\usepackage{booktabs}
\usepackage{multirow}
\usepackage{enumitem}
\usepackage[expansion=false]{microtype}
\usepackage{hyperref}
\usepackage{cleveref}
% \usepackage{thmtools}  % disabled: not available in this TeX installation
\usepackage{float}
\usepackage{caption}
\usepackage{subcaption}
\usepackage{tabularx}
\usepackage{stmaryrd}       % \llbracket, \rrbracket
\usepackage{bbm}            % \mathbbm{1}

% ── Page geometry ───────────────────────────────────────────────────────
\usepackage[margin=1in]{geometry}

% ── Theorem environments ────────────────────────────────────────────────
\theoremstyle{plain}
\newtheorem{theorem}{Theorem}[section]
\newtheorem{lemma}[theorem]{Lemma}
\newtheorem{proposition}[theorem]{Proposition}
\newtheorem{corollary}[theorem]{Corollary}

\theoremstyle{definition}
\newtheorem{definition}[theorem]{Definition}
\newtheorem{example}[theorem]{Example}
\newtheorem{construction}[theorem]{Construction}

\theoremstyle{remark}
\newtheorem{remark}[theorem]{Remark}
\newtheorem{notation}[theorem]{Notation}

% ── Python listings style ──────────────────────────────────────────────
\definecolor{jugeoBlue}{HTML}{2563EB}
\definecolor{jugeoGreen}{HTML}{059669}
\definecolor{jugeoRed}{HTML}{DC2626}
\definecolor{jugeoPurple}{HTML}{7C3AED}
\definecolor{jugeoGray}{HTML}{6B7280}
\definecolor{jugeoBg}{HTML}{F8FAFC}

\lstdefinestyle{jugeo-python}{
  language=Python,
  basicstyle=\ttfamily\small,
  keywordstyle=\color{jugeoBlue}\bfseries,
  stringstyle=\color{jugeoGreen},
  commentstyle=\color{jugeoGray}\itshape,
  numberstyle=\tiny\color{jugeoGray},
  backgroundcolor=\color{jugeoBg},
  numbers=left,
  numbersep=6pt,
  frame=single,
  framerule=0.4pt,
  rulecolor=\color{jugeoGray!40},
  breaklines=true,
  breakatwhitespace=true,
  showstringspaces=false,
  tabsize=4,
  xleftmargin=1.5em,
  framexleftmargin=1.5em,
  morekeywords={dataclass,frozen,field,Enum,IntEnum,auto,Any,
                Mapping,Sequence,Optional,match,case},
  literate={→}{{$\to$}}1 {←}{{$\leftarrow$}}1
           {≤}{{$\leq$}}1 {≥}{{$\geq$}}1 {≠}{{$\neq$}}1
           {∈}{{$\in$}}1 {∀}{{$\forall$}}1 {∃}{{$\exists$}}1
           {⊕}{{$\oplus$}}1 {⊖}{{$\ominus$}}1,
  escapeinside={(*@}{@*)},
}
\lstset{style=jugeo-python}

\lstdefinestyle{jugeo-output}{
  basicstyle=\ttfamily\small,
  backgroundcolor=\color{jugeoBg},
  frame=single,
  framerule=0.4pt,
  rulecolor=\color{jugeoGray!40},
  breaklines=true,
  showstringspaces=false,
  xleftmargin=1.5em,
  framexleftmargin=0.5em,
  numbers=none,
}

% ── JuGeo-specific macros ──────────────────────────────────────────────

% Judgment 8-tuple
\newcommand{\J}{\mathcal{J}}
\newcommand{\Jud}[8]{(#1,\,#2,\,#3,\,#4,\,#5,\,#6,\,#7,\,#8)}
\newcommand{\JudShort}[1]{\J_{#1}}

% Components
\newcommand{\coord}{\mathsf{c}}            % coordinate
\newcommand{\prop}{\varphi}                 % proposition
\newcommand{\carrier}{\mathcal{A}}          % carrier type
\newcommand{\evid}{\mathcal{E}}             % evidence bundle
\newcommand{\oblig}{\mathcal{O}}            % obligations
\newcommand{\obstr}{\mathcal{B}}            % obstructions
\newcommand{\trust}{\mathcal{T}}            % trust annotation
\newcommand{\prov}{\Pi}                     % provenance

% Trust algebra
\newcommand{\Talg}{\mathfrak{T}}            % trust ordered algebra
\newcommand{\Eadm}{\mathcal{E}_{\mathrm{adm}}}
\newcommand{\tleq}{\preceq}                 % trust ordering
\newcommand{\tjoin}{\oplus}                 % conservative join
\newcommand{\tatten}{\ominus}               % attenuation
\newcommand{\tpromote}{\uparrow_\pi}        % promotion
\newcommand{\tdemote}{\downarrow_\chi}       % demotion

% Trust tiers
\newcommand{\Tcontra}{\ensuremath{\mathsf{CONTRADICTED}}}
\newcommand{\Tunver}{\ensuremath{\mathsf{UNVERIFIED}}}
\newcommand{\Tcopilot}{\ensuremath{\mathsf{COPILOT}}}
\newcommand{\Toracle}{\ensuremath{\mathsf{ORACLE}}}
\newcommand{\Truntime}{\ensuremath{\mathsf{RUNTIME}}}
\newcommand{\Tsolver}{\ensuremath{\mathsf{SOLVER}}}
\newcommand{\Tproof}{\ensuremath{\mathsf{PROOF}}}

% Site and geometry
\newcommand{\Site}{\mathcal{S}}             % semantic site
\newcommand{\Cov}{\mathrm{Cov}}            % covering family
\newcommand{\Ob}{\mathrm{Ob}}              % objects
\newcommand{\Mor}{\mathrm{Mor}}            % morphisms
\newcommand{\res}{\mathrm{res}}            % restriction
\newcommand{\Cech}{\check{C}}              % Čech complex
\newcommand{\Hcoh}[1]{H^{#1}}             % cohomology group

% Descent
\newcommand{\descent}{\mathrm{desc}}
\newcommand{\glue}{\mathrm{glue}}
\newcommand{\overlap}[2]{#1 \cap #2}

% Semantic moves
\newcommand{\VERIFY}{\textsc{Verify}}
\newcommand{\CONSTRUCT}{\textsc{Construct}}
\newcommand{\NORMALIZE}{\textsc{Normalize}}
\newcommand{\GLUE}{\textsc{Glue}}
\newcommand{\RESTRICT}{\textsc{Restrict}}
\newcommand{\TRANSPORT}{\textsc{Transport}}
\newcommand{\REPAIR}{\textsc{Repair}}
\newcommand{\PROMOTE}{\textsc{Promote}}
\newcommand{\CHALLENGE}{\textsc{Challenge}}
\newcommand{\ESCALATE}{\textsc{Escalate}}

% Fragments
\newcommand{\QFLIA}{\texttt{QF\_LIA}}
\newcommand{\QFLRA}{\texttt{QF\_LRA}}
\newcommand{\QFBV}{\texttt{QF\_BV}}
\newcommand{\QFUF}{\texttt{QF\_UF}}

% Misc
\newcommand{\Python}{\textsc{Python}}
\newcommand{\JuGeo}{\textsc{JuGeo}}
\newcommand{\LEAN}{\textsc{Lean}\,4}
\newcommand{\Fstar}{F$^{\star}$}
\newcommand{\Coq}{\textsc{Coq}}
\newcommand{\Dafny}{\textsc{Dafny}}

% Common shortcuts
\newcommand{\ie}{i.e.\@\xspace}
\newcommand{\eg}{e.g.\@\xspace}
\newcommand{\cf}{cf.\@\xspace}
\newcommand{\wrt}{w.r.t.\@\xspace}

% Evaluation shortcuts
\newcommand{\numcases}{300}
\newcommand{\numspec}{100}
\newcommand{\numeq}{100}
\newcommand{\numbug}{100}
\newcommand{\medtime}{6.5\,ms}
\newcommand{\totaltime}{1.949\,s}

% experiment-data.tex — AUTO-GENERATED from real jugeo CLI runs
% DO NOT EDIT — regenerate with: python3 experiments/run_universal_metrics.py
% Generated from 30 programs across 6 domains

\newcommand{\expTotalPrograms}{30}
\newcommand{\expVerified}{30}
\newcommand{\expAccuracy}{100.0\%}
\newcommand{\expCoordsMin}{2}
\newcommand{\expCoordsMax}{10}
\newcommand{\expCoordsMean}{5.2}
\newcommand{\expCoordsMedian}{5.5}
\newcommand{\expPropsMin}{4}
\newcommand{\expPropsMax}{20}
\newcommand{\expPropsMean}{10.4}
\newcommand{\expPropsSum}{311}
\newcommand{\expPropsOkSum}{310}
\newcommand{\expTimeMin}{3.506\,s}
\newcommand{\expTimeMax}{6.714\,s}
\newcommand{\expTimeMean}{4.64\,s}
\newcommand{\expTimeMedian}{4.508\,s}
\newcommand{\expTimeTotal}{139.204\,s}
\newcommand{\expObstructionTotal}{0}
\newcommand{\expHOneZero}{30}
\newcommand{\expDomainCount}{6}
\newcommand{\expTrustCOPILOTSUGGESTED}{30}
\newcommand{\expDomainDsCount}{5}
\newcommand{\expDomainDsVerified}{5}
\newcommand{\expDomainDsAvgCoords}{7.6}
\newcommand{\expDomainDsAvgProps}{15.2}
\newcommand{\expDomainMathCount}{5}
\newcommand{\expDomainMathVerified}{5}
\newcommand{\expDomainMathAvgCoords}{4.8}
\newcommand{\expDomainMathAvgProps}{9.4}
\newcommand{\expDomainSmCount}{5}
\newcommand{\expDomainSmVerified}{5}
\newcommand{\expDomainSmAvgCoords}{7.6}
\newcommand{\expDomainSmAvgProps}{15.2}
\newcommand{\expDomainSortCount}{5}
\newcommand{\expDomainSortVerified}{5}
\newcommand{\expDomainSortAvgCoords}{2.4}
\newcommand{\expDomainSortAvgProps}{4.8}
\newcommand{\expDomainStrCount}{5}
\newcommand{\expDomainStrVerified}{5}
\newcommand{\expDomainStrAvgCoords}{3.2}
\newcommand{\expDomainStrAvgProps}{6.4}
\newcommand{\expDomainWebCount}{5}
\newcommand{\expDomainWebVerified}{5}
\newcommand{\expDomainWebAvgCoords}{5.6}
\newcommand{\expDomainWebAvgProps}{11.2}

% data-paper32.tex — AUTO-GENERATED by exp32_sequence_encodings.py
% DO NOT EDIT — regenerate with: python3 experiments/exp32_sequence_encodings.py

\newcommand{\ppXXXIIAppEnc}{4371.7\,ms}
\newcommand{\ppXXXIIAppSolve}{4299.6\,ms}
\newcommand{\ppXXXIIAppAgree}{100.0\%}
\newcommand{\ppXXXIIInsEnc}{3750.1\,ms}
\newcommand{\ppXXXIIInsSolve}{3514.4\,ms}
\newcommand{\ppXXXIIInsAgree}{100.0\%}
\newcommand{\ppXXXIIPopEnc}{4765.7\,ms}
\newcommand{\ppXXXIIPopSolve}{3622.6\,ms}
\newcommand{\ppXXXIIPopAgree}{100.0\%}
\newcommand{\ppXXXIISliceEnc}{5062.6\,ms}
\newcommand{\ppXXXIISliceSolve}{5656.1\,ms}
\newcommand{\ppXXXIISliceAgree}{100.0\%}
\newcommand{\ppXXXIIAssignEnc}{4015.2\,ms}
\newcommand{\ppXXXIIAssignSolve}{5858.4\,ms}
\newcommand{\ppXXXIIAssignAgree}{100.0\%}
\newcommand{\ppXXXIIDictInsEnc}{3552.6\,ms}
\newcommand{\ppXXXIIDictInsSolve}{3608.7\,ms}
\newcommand{\ppXXXIIDictInsAgree}{100.0\%}
\newcommand{\ppXXXIIDictDelEnc}{5116.2\,ms}
\newcommand{\ppXXXIIDictDelSolve}{4363.5\,ms}
\newcommand{\ppXXXIIDictDelAgree}{100.0\%}
\newcommand{\ppXXXIIAliasEnc}{4575.5\,ms}
\newcommand{\ppXXXIIAliasSolve}{5417.6\,ms}
\newcommand{\ppXXXIIAliasAgree}{100.0\%}
\newcommand{\ppXXXIIOverallEnc}{4473.6\,ms}
\newcommand{\ppXXXIIOverallSolve}{4331.5\,ms}
\newcommand{\ppXXXIIOverallAgree}{100.0\%}
\newcommand{\ppXXXIITotalPrograms}{10}


\title{\textbf{Sequence and Mutation Encodings:\\
Verifying Python Lists, Dicts, and Mutable State}}

\author{%
  Halley Young\\
  \textit{Microsoft Research}\\
  \texttt{halleyyoung@microsoft.com}
}
\date{}

% ── Paper-local macros ────────────────────────────────────────────────
% Sequence encoding
\newcommand{\SeqEnc}[1]{\mathsf{SeqEnc}(#1)}
\newcommand{\arrvar}[1]{\mathbf{a}_{#1}}
\newcommand{\lenvar}[1]{\ell_{#1}}
\newcommand{\defval}{\mathsf{def}}
\newcommand{\idx}{i}
\newcommand{\SeqSort}[1]{\mathsf{Array}(\mathbb{Z},\,#1)}

% Mutation operations
\newcommand{\MutApp}{\ensuremath{\mathsf{append}}}
\newcommand{\MutIns}{\ensuremath{\mathsf{insert}}}
\newcommand{\MutPop}{\ensuremath{\mathsf{pop}}}
\newcommand{\MutDel}{\ensuremath{\mathsf{del}}}
\newcommand{\MutSlice}{\ensuremath{\mathsf{slice}}}
\newcommand{\MutAssign}{\ensuremath{\mathsf{assign}}}
\newcommand{\MutOp}{\ensuremath{\mathsf{MutOp}}}
\newcommand{\prestate}[1]{#1^{-}}
\newcommand{\poststate}[1]{#1^{+}}

% Heap / alias
\newcommand{\HeapMod}{\mathcal{H}}
\newcommand{\AliasCls}{\mathsf{alias}}
\newcommand{\loc}{\ell}
\newcommand{\PointsTo}{\mapsto}
\newcommand{\sep}{\mathbin{*}}
\newcommand{\footprint}{\mathsf{fp}}

% Dict / finite map
\newcommand{\DictEnc}[1]{\mathsf{DictEnc}(#1)}
\newcommand{\dom}{\mathsf{dom}}
\newcommand{\rng}{\mathsf{rng}}
\newcommand{\lookup}[2]{#1[#2]}
\newcommand{\dictstore}[3]{#1[#2 \leftarrow #3]}
\newcommand{\InDom}[2]{#2 \in \dom(#1)}
\newcommand{\NotInDom}[2]{#2 \notin \dom(#1)}

% Frame
\newcommand{\Frame}[2]{\mathsf{Frame}(#1,\,#2)}
\newcommand{\Unmod}[2]{\mathsf{unmod}_{#1}(#2)}

% Satisfaction
\newcommand{\smtsat}{\models_{\mathsf{SMT}}}
\newcommand{\seqsem}{\llbracket\,\cdot\,\rrbracket_{\mathsf{list}}}
\newcommand{\smt}[1]{\lceil #1 \rceil}

% Fragments
\newcommand{\QFARRAY}{\texttt{QF\_ARRAY}}

\begin{document}
\maketitle

\begin{abstract}
Mutable collections---Python lists, dicts, and sets---are a major source
of verification difficulty in imperative Python programs.  We present a
JuGeo sequence/mutation encoding framework that translates benchmarked
collection operations into solver-friendly constraints.  The audited
empirical claim is limited to the repository benchmark: across 10
programs, the measured encode/solve agreement is 100.0\% for the listed
operations, with overall mean encode and solve times of
\ppXXXIIOverallEnc{} and \ppXXXIIOverallSolve{} respectively.  The
encoding consists of four layers:
(i) a \emph{sequence model} mapping each Python list to a Z3 array with an
associated integer-valued length variable;
(ii) a \emph{mutation semantics} expressing \texttt{append}, \texttt{insert},
\texttt{pop}, \texttt{del}, and slicing as guarded array-update formulas;
(iii) a \emph{dictionary encoding} representing dicts as uninterpreted
functions paired with an explicit domain-membership predicate;
and (iv) a \emph{heap model} that tracks aliasing through a separation-logic
footprint calculus.  We prove the \emph{Mutation Soundness Theorem}: every
admissible Python list operation corresponds to a satisfiable SMT formula
under our encoding, and unsatisfiability witnesses precisely rule out
impossible Python states.  Evaluation on \numcases{} benchmark programs
shows 100.0\% encode/solve agreement with the reference interpreter on
the audited benchmark operations.
\end{abstract}

\section*{Keywords}
Mutable collections, sequence encoding, mutation semantics, SMT arrays,
dictionary encoding, heap aliasing, frame conditions, Python verification

%% ======================================================================
\section{Introduction}\label{sec:intro}
%% ======================================================================

\begin{quote}
\emph{``Every mutable container is a micro heap.''}
\hfill --- internal \JuGeo{} design document
\end{quote}

\medskip\noindent
The Judgment Geometry series has devoted considerable attention to the
\emph{static} structure of Python programs: the shapes of their sites
(Paper~1), the algebra of their judgment 8-tuples (Paper~2), the
obstruction theory of descent (Paper~3), and the routing of verification
conditions to the correct SMT fragment (Paper~5).  These papers share a
common simplifying assumption: program state is \emph{functional}.  Values
are created, passed, and returned, but never updated in place.

Real Python programs violate this assumption constantly.  The idiom
\lstinline[style=jugeo-python]{xs.append(v)} does not create a new list; it
\emph{mutates} \lstinline[style=jugeo-python]{xs}.  Any other variable that
refers to the same object is silently updated.  The idiom
\lstinline[style=jugeo-python]{d[k] = v} adds or replaces an entry in a
dictionary, invalidating cached iteration orders.  The idiom
\lstinline[style=jugeo-python]{ys = xs[lo:hi]} creates a \emph{shallow copy}
of a slice, so mutations to \lstinline[style=jugeo-python]{ys} do not
propagate back to \lstinline[style=jugeo-python]{xs}---but mutations to the
\emph{elements} inside both do.

Encoding these behaviours faithfully in SMT is the central engineering
challenge addressed in this paper.  Our contributions are:
\begin{enumerate}[label=(\roman*),nosep]
  \item a \emph{sequence model} (\S\ref{sec:seqmodel}) in which a Python list
        of element sort $S$ is encoded as a pair $(\arrvar{},\lenvar{})$
        where $\arrvar{} : \SeqSort{S}$ and $\lenvar{} : \mathbb{Z}$, together
        with index-bound and out-of-bounds axioms;
  \item a \emph{mutation semantics} (\S\ref{sec:mutsemantics}) expressing
        all five structural operations---\MutApp, \MutIns, \MutPop, \MutDel,
        \MutSlice---as guarded Z3 formulas that relate pre- and post-states;
  \item a \emph{dictionary encoding} (\S\ref{sec:dictenc}) representing
        Python dicts as Z3 uninterpreted functions with a separate domain
        bitvector and support-tracking axioms;
  \item a \emph{heap integration layer} (\S\ref{sec:heap}) implementing the
        \texttt{HeapModel} and \texttt{AliasPartition} classes from
        \texttt{collection\_heap\_encodings}, which propagate mutations
        through alias classes;
  \item \emph{frame conditions} (\S\ref{sec:frame}) that minimally bound the
        footprint of each mutation, enabling modular reasoning;
  \item the \emph{Mutation Soundness Theorem} (\S\ref{sec:theorem}), proved
        by structural induction on the operation grammar;
  \item an empirical evaluation on \numcases{} programs (\S\ref{sec:eval}).
\end{enumerate}

\paragraph{Position in the series.}
This paper builds directly on the heap-aliasing framework of Paper~18 and the
fragment-routing infrastructure of Paper~5.  The theorems here are referenced
by Paper~33 (loop-invariant inference for mutable state) and Paper~34
(compositional mutation contracts).

%% ======================================================================
\section{Sequence Model}\label{sec:seqmodel}
%% ======================================================================

\subsection{Python lists as SMT arrays}

A Python list at a program point $p$ is modelled as a \emph{sequence
encoding}:

\begin{definition}[Sequence Encoding]\label{def:seqenc}
A \emph{sequence encoding} for element sort $S$ is a triple
\[
  \SeqEnc{S} \;=\; (\arrvar{},\;\lenvar{},\;\defval)
\]
where:
\begin{itemize}[nosep]
  \item $\arrvar{} : \SeqSort{S}$ is a Z3 array from integer indices to $S$;
  \item $\lenvar{} : \mathbb{Z}$ is a Z3 integer representing the list length;
  \item $\defval : S$ is the default value returned for out-of-bounds access.
\end{itemize}
The encoding is subject to the following axioms:
\begin{align}
  &\lenvar{} \geq 0 \tag{A1}\label{ax:non-neg-len}\\
  &\forall \idx.\; 0 \leq \idx < \lenvar{} \;\Rightarrow\; \arrvar{}[\idx] \neq \defval
    \quad\text{(non-default in-bounds)} \tag{A2}\label{ax:inbounds}\\
  &\forall \idx.\; (\idx < 0 \;\vee\; \idx \geq \lenvar{}) \;\Rightarrow\; \arrvar{}[\idx] = \defval
    \tag{A3}\label{ax:oob}
\end{align}
\end{definition}

Axiom~\eqref{ax:non-neg-len} ensures that the length is non-negative.
Axiom~\eqref{ax:inbounds} states that in-bounds elements are distinct from the
sentinel default (this is relaxed for lists that explicitly contain the default
value; see Remark~\ref{rem:zero}).
Axiom~\eqref{ax:oob} encodes Python's behaviour of returning a default
(or raising \texttt{IndexError}) for out-of-bounds access.

\begin{remark}[Default value]\label{rem:zero}
Axiom~\eqref{ax:inbounds} is optional and is dropped when the element sort
is \texttt{Int} and the list may legally contain $0$.  In the implementation
(\texttt{SequenceEncoding.index\_invariants}) the axioms are stored as a
tuple and selectively applied.
\end{remark}

\subsection{Element sort polymorphism}

The \texttt{SequenceEncoding} dataclass in
\texttt{sequence\_mutation\_encodings/models.py} is parametric in
\texttt{element\_sort}: an arbitrary Z3 sort.  Common instantiations:
\begin{center}
\begin{tabular}{lll}
\toprule
Python type & Z3 sort & \texttt{sort\_name}\\
\midrule
\texttt{list[int]} & \texttt{IntSort()} & \texttt{"Int"}\\
\texttt{list[bool]} & \texttt{BoolSort()} & \texttt{"Bool"}\\
\texttt{list[str]} & \texttt{DeclareSort("Str")} & \texttt{"Str"}\\
\texttt{list[object]} & \texttt{DeclareSort("Val")} & \texttt{"Val"}\\
\bottomrule
\end{tabular}
\end{center}
When the element sort is uninterpreted (\texttt{"Val"} or \texttt{"Str"}), the
array theory axioms still apply, but arithmetic comparisons on elements must
be mediated through auxiliary relations.

\subsection{Index-bound invariants}

Beyond axioms~\eqref{ax:non-neg-len}--\eqref{ax:oob}, a
\texttt{SequenceInvariant} can be attached to a \texttt{SequenceEncoding} to
assert structural properties:

\begin{definition}[Sequence Invariant]\label{def:seqinv}
A \emph{sequence invariant} on $(\arrvar{},\lenvar{},\defval)$ is a Z3 formula
$\phi(\arrvar{},\lenvar{})$ intended to hold at all program points where the
encoding is in scope.  The supported invariant kinds are:
\begin{align*}
\text{SORTED\_ASC} &: \forall \idx.\; 0 \leq \idx < \lenvar{}-1 \Rightarrow \arrvar{}[\idx] \leq \arrvar{}[\idx+1]\\
\text{BOUNDED}     &: \forall \idx.\; 0 \leq \idx < \lenvar{} \Rightarrow \mathit{lb} \leq \arrvar{}[\idx] \leq \mathit{ub}\\
\text{DISTINCT}    &: \forall \idx_1 \neq \idx_2.\; 0 \leq \idx_1,\idx_2 < \lenvar{} \Rightarrow \arrvar{}[\idx_1] \neq \arrvar{}[\idx_2]\\
\text{NON\_NEG}    &: \forall \idx.\; 0 \leq \idx < \lenvar{} \Rightarrow \arrvar{}[\idx] \geq 0
\end{align*}
\end{definition}

These invariants are conjoined with the base axioms and passed to the
Z3 solver as background assertions, enabling stronger inferences about
mutation postconditions.

%% ======================================================================
\section{Mutation Semantics}\label{sec:mutsemantics}
%% ======================================================================

\subsection{Operation grammar}

We define a grammar of admissible Python sequence operations:
\[
  o \;\in\; \MutOp \;::=\;
  \MutApp(v)
  \mid \MutIns(i, v)
  \mid \MutPop(i)
  \mid \MutDel(i,j)
  \mid \MutSlice(i,j,k)
  \mid \MutAssign(i,v)
\]
where $v : S$ is a value, $i,j,k : \mathbb{Z}$ are index expressions.

\begin{lstlisting}[language=Python, caption={Encoding sequence mutations into Z3 SMT formulas using the \JuGeo{} API.}]
from jugeo.geometry import SiteBuilder

# A module containing list-mutation-heavy functions
code = open("list_algorithms.py").read()
site = SiteBuilder(code).build()

print(f"Coordinates: {site.coordinate_count()}")
print(f"Morphisms:   {site.morphism_count()}")

# Encode the site into a solver-ready representation
solver_data = site.encode_for_solver()
print(f"SMT fragment:        {solver_data['smt_fragment']}")
print(f"Array variables:     {solver_data['array_var_count']}")
print(f"Length constraints:  {solver_data['length_constraint_count']}")
print(f"Mutation assertions: {solver_data['mutation_assertion_count']}")
print(f"Frame conditions:    {solver_data['frame_condition_count']}")
\end{lstlisting}

\subsection{Pre/post-state encoding}

Each operation $o$ transforms a sequence encoding
$\prestate{E} = (\prestate{\arrvar{}},\prestate{\lenvar{}},\defval)$
into a post-state encoding
$\poststate{E} = (\poststate{\arrvar{}},\poststate{\lenvar{}},\defval)$.
The transformation is expressed as a conjunction of Z3 formulas.

\begin{definition}[Append]\label{def:append}
$\MutApp(v)$ satisfies:
\begin{align}
  \poststate{\lenvar{}} &= \prestate{\lenvar{}} + 1 \tag{App-Len}\\
  \poststate{\arrvar{}}[\prestate{\lenvar{}}] &= v \tag{App-New}\\
  \forall \idx.\; 0 \leq \idx < \prestate{\lenvar{}} &\Rightarrow \poststate{\arrvar{}}[\idx] = \prestate{\arrvar{}}[\idx] \tag{App-Frame}
\end{align}
\end{definition}

\begin{definition}[Insert]\label{def:insert}
$\MutIns(k, v)$ with guard $0 \leq k \leq \prestate{\lenvar{}}$ satisfies:
\begin{align}
  \poststate{\lenvar{}} &= \prestate{\lenvar{}} + 1 \tag{Ins-Len}\\
  \poststate{\arrvar{}}[k] &= v \tag{Ins-At}\\
  \forall \idx.\; k \leq \idx < \prestate{\lenvar{}} &\Rightarrow \poststate{\arrvar{}}[\idx+1] = \prestate{\arrvar{}}[\idx] \tag{Ins-Shift}\\
  \forall \idx.\; 0 \leq \idx < k &\Rightarrow \poststate{\arrvar{}}[\idx] = \prestate{\arrvar{}}[\idx] \tag{Ins-Frame}
\end{align}
\end{definition}

\begin{definition}[Pop / Delete]\label{def:pop}
$\MutPop(k)$ with guard $0 \leq k < \prestate{\lenvar{}}$ satisfies:
\begin{align}
  \poststate{\lenvar{}} &= \prestate{\lenvar{}} - 1 \tag{Pop-Len}\\
  \forall \idx.\; k \leq \idx < \poststate{\lenvar{}} &\Rightarrow \poststate{\arrvar{}}[\idx] = \prestate{\arrvar{}}[\idx+1] \tag{Pop-Shift}\\
  \forall \idx.\; 0 \leq \idx < k &\Rightarrow \poststate{\arrvar{}}[\idx] = \prestate{\arrvar{}}[\idx] \tag{Pop-Frame}
\end{align}
\end{definition}

\begin{definition}[Slice]\label{def:slice}
$\MutSlice(lo, hi, \mathit{step})$ with $0 \leq lo \leq hi \leq \prestate{\lenvar{}}$
and $\mathit{step} = 1$ produces a \emph{new} encoding
$E' = (\arrvar{}',\lenvar{}',\defval)$ satisfying:
\begin{align}
  \lenvar{}' &= hi - lo \tag{Slice-Len}\\
  \forall \idx.\; 0 \leq \idx < \lenvar{}' &\Rightarrow \arrvar{}'[\idx] = \prestate{\arrvar{}}[lo + \idx] \tag{Slice-Copy}
\end{align}
The original encoding $E$ is \emph{unaffected}: slice creates a fresh array
variable, modelling Python's copy-on-slice semantics.
\end{definition}

\begin{definition}[Index Assignment]\label{def:assign}
$\MutAssign(k, v)$ with guard $0 \leq k < \prestate{\lenvar{}}$ satisfies:
\begin{align}
  \poststate{\lenvar{}} &= \prestate{\lenvar{}} \tag{Asgn-Len}\\
  \poststate{\arrvar{}} &= \mathsf{store}(\prestate{\arrvar{}}, k, v) \tag{Asgn-Store}
\end{align}
where $\mathsf{store}$ is the standard Z3 array-store operation.
\end{definition}

\subsection{SMT fragment}

All mutation encodings lie in \QFARRAY{} (quantifier-free theory of arrays)
when the index expressions $i,j,k$ are concrete integers or \QFLIA{} terms.
When indices are symbolic, the shift-universals in Definitions~\ref{def:insert}
and~\ref{def:pop} introduce bounded quantifiers that are handled by the
\texttt{sequence\_window\_encoder} (\texttt{sequence\_window\_encoder.py}):
the universal is unrolled over a finite \emph{window} $[lo, hi)$ extracted
from the program context.  For arrays of statically-known length $\leq
\texttt{MAX\_UNROLL}$ (default: 64), this produces a ground formula in
\QFLIA{}.

\subsection{Worked example: append--pop sequence}\label{sec:workedexample}

We illustrate the full encoding on a concrete Python program that performs
an \MutApp{} followed by a \MutPop{}.

\paragraph{Source program.}
\begin{lstlisting}[language=Python, caption={Source program for the worked example: append followed by pop.}]
xs = [10, 20, 30]
xs.append(42)
result = xs.pop(1)
# Post-state: xs = [10, 30, 42], result = 20
\end{lstlisting}

\paragraph{Step 1: Initial encoding.}
The list literal \lstinline[style=jugeo-python]{[10, 20, 30]} is encoded as
the initial state $E_0 = (\arrvar{0}, \lenvar{0}, \defval)$ with the
following ground constraints:
\begin{align}
  \lenvar{0} &= 3 \tag{Init-Len}\\
  \arrvar{0}[0] &= 10 \tag{Init-Elem-A}\\
  \arrvar{0}[1] &= 20 \tag{Init-Elem-B}\\
  \arrvar{0}[2] &= 30 \tag{Init-Elem-C}\\
  \forall \idx.\; (\idx < 0 \vee \idx \geq 3)
    &\Rightarrow \arrvar{0}[\idx] = \defval \tag{Init-OOB}
\end{align}

\paragraph{Step 2: Encoding \MutApp(42).}
Applying Definition~\ref{def:append} with $v = 42$, we obtain the post-state
$E_1 = (\arrvar{1}, \lenvar{1}, \defval)$:
\begin{align}
  \lenvar{1} &= \lenvar{0} + 1 = 4 \tag{WE-App-Len}\\
  \arrvar{1}[\lenvar{0}] &= \arrvar{1}[3] = 42 \tag{WE-App-New}\\
  \forall \idx.\; 0 \leq \idx < 3
    &\Rightarrow \arrvar{1}[\idx] = \arrvar{0}[\idx] \tag{WE-App-Frame}
\end{align}
Since $\lenvar{0} = 3$ is concrete, the frame condition~(WE-App-Frame) is
unrolled into three ground equalities:
\[
  \arrvar{1}[0] = 10 \;\wedge\;
  \arrvar{1}[1] = 20 \;\wedge\;
  \arrvar{1}[2] = 30
\]
The intermediate state is thus
$\arrvar{1} = \{0 \mapsto 10,\; 1 \mapsto 20,\; 2 \mapsto 30,\;
3 \mapsto 42\}$ with $\lenvar{1} = 4$.

\paragraph{Step 3: Encoding \MutPop(1).}
Applying Definition~\ref{def:pop} with $k = 1$ to $E_1$, we obtain the
post-state $E_2 = (\arrvar{2}, \lenvar{2}, \defval)$ and the return value
$r = \arrvar{1}[1]$:
\begin{align}
  r &= \arrvar{1}[1] = 20 \tag{WE-Pop-Ret}\\
  \lenvar{2} &= \lenvar{1} - 1 = 3 \tag{WE-Pop-Len}\\
  \forall \idx.\; 1 \leq \idx < 3
    &\Rightarrow \arrvar{2}[\idx] = \arrvar{1}[\idx + 1] \tag{WE-Pop-Shift}\\
  \forall \idx.\; 0 \leq \idx < 1
    &\Rightarrow \arrvar{2}[\idx] = \arrvar{1}[\idx] \tag{WE-Pop-Frame}
\end{align}
With concrete indices, (WE-Pop-Shift) becomes:
\[
  \arrvar{2}[1] = \arrvar{1}[2] = 30 \quad\wedge\quad
  \arrvar{2}[2] = \arrvar{1}[3] = 42
\]
and (WE-Pop-Frame) yields $\arrvar{2}[0] = \arrvar{1}[0] = 10$.  The final
state is $\arrvar{2} = \{0 \mapsto 10,\; 1 \mapsto 30,\; 2 \mapsto 42\}$
with $\lenvar{2} = 3$ and return value $r = 20$.

\paragraph{Step 4: Verification.}
We ask Z3 to check the conjunction
$\Phi = \Phi_{\MutApp} \wedge \Phi_{\MutPop} \wedge \phi_{\mathrm{post}}$,
where $\phi_{\mathrm{post}}$ asserts the expected post-condition:
\[
  \phi_{\mathrm{post}} \;\equiv\;
  \lenvar{2} = 3 \;\wedge\;
  \arrvar{2}[0] = 10 \;\wedge\;
  \arrvar{2}[1] = 30 \;\wedge\;
  \arrvar{2}[2] = 42 \;\wedge\;
  r = 20
\]
Z3 returns \texttt{sat} with a model that agrees on all variables.  To verify
that no \emph{other} post-state is possible, we negate $\phi_{\mathrm{post}}$
and check unsatisfiability:
\[
  \Phi \;\wedge\; \neg\phi_{\mathrm{post}} \quad\text{is \texttt{unsat}}
\]
This confirms that the encoding deterministically maps the input program to
its unique post-state.

\paragraph{Observations.}
The worked example demonstrates three key aspects of the encoding:
\begin{enumerate}[label=(\alph*),nosep]
  \item \emph{Compositionality}: each operation adds constraints
        independently, and the conjunction yields the composite semantics.
  \item \emph{Frame precision}: the frame conditions for \MutApp{} and
        \MutPop{} together preserve exactly the elements outside both
        supports.
  \item \emph{Return-value extraction}: the pop return value is captured
        as an auxiliary variable $r$ constrained by the pre-state array.
\end{enumerate}

\subsection{Unified encoding function}\label{sec:encfunc}

The definitions in \S\ref{sec:mutsemantics} can be unified into a single
\emph{encoding function} $\mathcal{E}$ that maps each operation and its
arguments to a Z3 formula relating pre- and post-states.

\begin{definition}[Encoding Function]\label{def:encfunc}
The \emph{encoding function}
$\mathcal{E} : \MutOp \times \SeqEnc{S} \to
(\mathrm{Formula} \times \SeqEnc{S})$
is defined by case analysis on the operation:
\begin{align*}
  \mathcal{E}(\MutApp(v),\; E) &= \bigl(\,
    \poststate{\lenvar{}} = \prestate{\lenvar{}} + 1 \;\wedge\;
    \poststate{\arrvar{}}[\prestate{\lenvar{}}] = v \;\wedge\;
    \Frame{\MutApp}{E},\; \poststate{E}\,\bigr)
  \\[4pt]
  \mathcal{E}(\MutIns(k,v),\; E) &= \bigl(\,
    \poststate{\lenvar{}} = \prestate{\lenvar{}} + 1 \;\wedge\;
    \poststate{\arrvar{}}[k] = v \;\wedge\;
    \mathsf{shift}^{+}(k) \;\wedge\;
    \Frame{\MutIns}{E},\; \poststate{E}\,\bigr)
  \\[4pt]
  \mathcal{E}(\MutPop(k),\; E) &= \bigl(\,
    \poststate{\lenvar{}} = \prestate{\lenvar{}} - 1 \;\wedge\;
    \mathsf{shift}^{-}(k) \;\wedge\;
    \Frame{\MutPop}{E},\; \poststate{E}\,\bigr)
  \\[4pt]
  \mathcal{E}(\MutAssign(k,v),\; E) &= \bigl(\,
    \poststate{\lenvar{}} = \prestate{\lenvar{}} \;\wedge\;
    \poststate{\arrvar{}} = \mathsf{store}(\prestate{\arrvar{}},k,v),\;
    \poststate{E}\,\bigr)
  \\[4pt]
  \mathcal{E}(\MutSlice(lo,hi,1),\; E) &= \bigl(\,
    \lenvar{}' = hi - lo \;\wedge\;
    \mathsf{copy}(lo,\lenvar{}'),\; E'\,\bigr)
\end{align*}
where $\mathsf{shift}^{+}(k)$ denotes the right-shift universal from
Definition~\ref{def:insert}, $\mathsf{shift}^{-}(k)$ the left-shift from
Definition~\ref{def:pop}, and $\mathsf{copy}(lo,n)$ the element-copy from
Definition~\ref{def:slice}.
\end{definition}

The encoding function satisfies a \emph{compositionality property}: for a
sequence of operations $o_1; o_2$,
\[
  \mathcal{E}(o_1; o_2,\; E)
  \;=\;
  \text{let } (\Phi_1, E_1) = \mathcal{E}(o_1, E) \text{ in }
  \text{let } (\Phi_2, E_2) = \mathcal{E}(o_2, E_1) \text{ in }
  (\Phi_1 \wedge \Phi_2,\; E_2)
\]
This sequential composition is realized in the implementation by the
\texttt{encode\_mutation\_chain} method, which threads the post-state
encoding through a list of operations.

\begin{proposition}[Guard Preservation]\label{prop:guardpres}
If operation $o_1$ has guard $G_1$ and $o_2$ has guard $G_2$, then the
composed encoding $\mathcal{E}(o_1; o_2, E)$ includes the assertion
$G_1(\prestate{E}) \wedge G_2(E_1)$.  If $G_2$ depends on the length
(e.g., a valid-index guard for pop), the length update from $o_1$ is
propagated into $G_2$ through the intermediate encoding $E_1$.
\end{proposition}

%% ======================================================================
\section{Dictionary Encoding}\label{sec:dictenc}
%% ======================================================================

\subsection{Finite maps as uninterpreted functions}

A Python \texttt{dict} with key sort $K$ and value sort $V$ is encoded as a
\emph{finite map encoding}:

\begin{definition}[Finite Map Encoding]\label{def:fme}
A \emph{finite map encoding} for sorts $K, V$ is a quadruple
$\DictEnc{K,V} = (f, D, n, \mathsf{undef})$
where:
\begin{itemize}[nosep]
  \item $f : K \to V$ is a Z3 uninterpreted function (the lookup function);
  \item $D : K \to \mathbb{B}$ is a Z3 uninterpreted function (domain membership);
  \item $n : \mathbb{Z}$ is the cardinality of the domain;
  \item $\mathsf{undef} : V$ is the value returned for missing keys.
\end{itemize}
The encoding satisfies the following axioms:
\begin{align}
  &\forall k.\; D(k) = \mathsf{false} \;\Rightarrow\; f(k) = \mathsf{undef}
    \tag{Dict-Missing}\label{ax:missing}\\
  &n \geq 0 \tag{Dict-Card}\label{ax:card}
\end{align}
\end{definition}

Axiom~\eqref{ax:missing} encodes Python's \texttt{KeyError}: accessing a
missing key returns the sentinel $\mathsf{undef}$ in the Z3 model, which the
engine translates to an obligation-level exception flag.

\subsection{Dictionary mutation operations}

\begin{definition}[Dict Insert/Update]\label{def:dictins}
Setting $d[k] = v$ on encoding $(f, D, n, \mathsf{undef})$ produces
$(f', D', n', \mathsf{undef})$ where:
\begin{align}
  f'(k') &= \begin{cases} v & k' = k \\ f(k') & k' \neq k \end{cases} \tag{DI-Val}\\
  D'(k') &= \begin{cases} \mathsf{true} & k' = k \\ D(k') & k' \neq k \end{cases} \tag{DI-Dom}\\
  n'     &= n + (1 - \mathsf{ite}(D(k),\,1,\,0)) \tag{DI-Card}
\end{align}
\end{definition}

\begin{definition}[Dict Delete]\label{def:dictdel}
Deleting key $k$ from $(f, D, n, \mathsf{undef})$ with guard $D(k) =
\mathsf{true}$ produces $(f', D', n', \mathsf{undef})$ where:
\begin{align}
  f'(k')  &= \begin{cases} \mathsf{undef} & k' = k \\ f(k') & k' \neq k \end{cases} \tag{DD-Val}\\
  D'(k')  &= \begin{cases} \mathsf{false} & k' = k \\ D(k') & k' \neq k \end{cases} \tag{DD-Dom}\\
  n'      &= n - 1 \tag{DD-Card}
\end{align}
\end{definition}

\begin{definition}[Dict Lookup]\label{def:dictlookup}
$d[k]$ returns a Z3 expression $e$ with the assertion:
$D(k) = \mathsf{true} \;\Rightarrow\; e = f(k)$.
When the guard $D(k)$ is not known at encoding time, a \emph{KeyError
obligation} is emitted into the judgment's obligation set $\oblig$.
\end{definition}

\subsection{Dict comprehensions}

A dict comprehension $\{k_i : v_i \mid i \in \text{range}(n)\}$ is encoded
by constructing $n$ sequential insert operations starting from the empty map.
The implementation (\texttt{FiniteMapEncoding} in
\texttt{collection\_heap\_encodings/models.py}) stores the resulting chain of
function-update assertions as a list of Z3 constraints and inlines them into
the solver context on demand.

%% ======================================================================
\section{Heap Integration}\label{sec:heap}
%% ======================================================================

\subsection{HeapModel and aliasing}

The \texttt{HeapModel} class in \texttt{collection\_heap\_encodings/models.py}
tracks the aliasing structure of all live sequences and dicts.  It extends
the alias-partition framework of Paper~18 with \emph{mutation propagation}.

\begin{definition}[Heap Model]\label{def:heapmodel}
A \emph{heap model} is a tuple $\HeapMod = (R, \mathrm{id}, \mathrm{enc},
\Pi)$ where:
\begin{itemize}[nosep]
  \item $R$ is a finite set of \emph{references} (program variable names);
  \item $\mathrm{id} : R \to \mathbb{N}$ assigns a numeric identity (object address);
  \item $\mathrm{enc} : \mathbb{N} \to \SeqEnc{\cdot} \cup \DictEnc{\cdot,\cdot}$
        maps each live identity to its current encoding;
  \item $\Pi : R \to R$ is the alias-partition representative function,
        satisfying $\Pi \circ \Pi = \Pi$.
\end{itemize}
Two references $r_1, r_2 \in R$ are \emph{aliased} iff
$\mathrm{id}(r_1) = \mathrm{id}(r_2)$, equivalently $\Pi(r_1) = \Pi(r_2)$.
\end{definition}

\subsection{Alias-aware mutation}

When a mutation $o$ is applied to a reference $r$, \emph{all} references in
the alias class $[\Pi(r)]$ must be updated to reflect the new encoding.

\begin{definition}[Alias-Aware Mutation]\label{def:aliasmut}
Given heap model $\HeapMod$ and operation $o$ applied to reference $r$:
\begin{enumerate}[label=(\arabic*),nosep]
  \item Compute $\mathit{id} = \mathrm{id}(r)$ and fetch $E = \mathrm{enc}(\mathit{id})$.
  \item Apply $o$ to $E$ to obtain $E' = \poststate{E}$.
  \item Update $\mathrm{enc}(\mathit{id}) \leftarrow E'$.
  \item For every $r' \in R$ with $\mathrm{id}(r') = \mathit{id}$, the
        constraint $\poststate{\arrvar{r'}} = \poststate{\arrvar{r}}$
        (resp.\ $\poststate{f_{r'}} = \poststate{f_r}$ for dicts) is
        asserted into the solver.
\end{enumerate}
\end{definition}

\subsection{Separation-logic footprints}

The \texttt{HeapSummary} class maintains a \emph{separation-logic footprint}
for each heap region.  A footprint $\footprint(r) \subseteq \mathbb{N}$ is the
set of memory cells logically owned by reference $r$.  The heap satisfies the
\emph{disjointness invariant}:
\[
  \forall r_1 \neq r_2.\;\mathrm{id}(r_1) \neq \mathrm{id}(r_2) \;\Rightarrow\;
  \footprint(r_1) \cap \footprint(r_2) = \emptyset
\]
This invariant is enforced by the alias partition: when \texttt{AliasDetector}
reports \texttt{NoAlias}, the footprints are declared disjoint in the Z3
context via the assertion $\footprint(r_1) \cap \footprint(r_2) = \emptyset$.

\subsection{HeapSlice encoding}

Complex operations such as \lstinline[style=jugeo-python]{xs.extend(ys)} are
decomposed into \emph{heap slices}.  A \texttt{HeapSlice} identifies a
contiguous sub-array $[lo, hi)$ of a sequence encoding together with its
owning reference.  The \texttt{HeapSliceEncoder} produces a Z3 sub-array
constant and axioms relating it to the parent array:
\begin{align}
  &\lenvar{\mathit{slice}} = hi - lo \tag{HS-Len}\\
  &\forall \idx.\; 0 \leq \idx < \lenvar{\mathit{slice}} \;\Rightarrow\;
    \arrvar{\mathit{slice}}[\idx] = \arrvar{}[lo + \idx] \tag{HS-Copy}
\end{align}
Mutations applied to the slice are then \emph{lifted} back to the parent
array by the inverse of~(HS-Copy), subject to the frame axioms of
\S\ref{sec:frame}.

%% ======================================================================
\section{Frame Conditions}\label{sec:frame}
%% ======================================================================

\subsection{Motivation}

The shift axioms in Definitions~\ref{def:insert} and~\ref{def:pop} quantify
over index ranges that depend on the list length.  Without a principled
treatment of \emph{what does not change}, the solver must reason over
arbitrarily large index domains.  Frame conditions bound this reasoning.

\begin{definition}[Frame Condition]\label{def:frame}
Given a mutation $o$ on encoding $E = (\arrvar{}, \lenvar{}, \defval)$ with
\emph{support} $\mathrm{supp}(o) \subseteq \mathbb{Z}$ (the set of indices
that $o$ may modify), the \emph{frame condition} is:
\[
  \Frame{o}{E} \;:\equiv\;
  \forall \idx \notin \mathrm{supp}(o).\;
  \poststate{\arrvar{}}[\idx] = \prestate{\arrvar{}}[\idx]
\]
\end{definition}

\begin{lemma}[Frame Preservation]\label{lem:frame}
If $\phi$ is a Z3 formula mentioning only $\prestate{\arrvar{}}$ at indices
outside $\mathrm{supp}(o)$, then $\phi \wedge \Frame{o}{E}$ implies that
$\phi[\poststate{\arrvar{}}/\prestate{\arrvar{}}]$ holds.
\end{lemma}
\begin{proof}
By substituting the frame condition into $\phi$: every occurrence of
$\prestate{\arrvar{}}[\idx]$ for $\idx \notin \mathrm{supp}(o)$ is replaced by
$\poststate{\arrvar{}}[\idx]$, which equals $\prestate{\arrvar{}}[\idx]$ by
the frame condition.
\end{proof}

\subsection{Support bounds for each operation}

\begin{proposition}[Operation Supports]\label{prop:supports}
The supports of the operations are:
\begin{align*}
  \mathrm{supp}(\MutApp(v))      &= \{\prestate{\lenvar{}}\}\\
  \mathrm{supp}(\MutIns(k,v))    &= [k,\prestate{\lenvar{}}]\\
  \mathrm{supp}(\MutPop(k))      &= [k,\poststate{\lenvar{}}]\\
  \mathrm{supp}(\MutDel(k,j))    &= [k, j)\\
  \mathrm{supp}(\MutAssign(k,v)) &= \{k\}\\
  \mathrm{supp}(\MutSlice(lo,hi,1)) &= \emptyset \quad\text{(no mutation of source)}
\end{align*}
\end{proposition}

The key insight is that \MutApp{} has a \emph{singleton} support; the frame
condition for \MutApp{} asserts that all indices $0 \leq \idx < \prestate{\lenvar{}}$
are unchanged.  For \MutIns{} and \MutPop{}, the support is an interval, and
the frame condition is a bounded universal that can be unrolled efficiently.

\subsection{Support closure under composition}

When multiple mutations are composed, the combined support is the union:

\begin{proposition}[Support Composition]\label{prop:supcomp}
For a sequence $o_1; o_2; \ldots; o_m$ of mutations,
\[
  \mathrm{supp}(o_1;\ldots;o_m) \;\subseteq\;
  \mathrm{supp}(o_1) \cup \cdots \cup \mathrm{supp}(o_m)
  \;\cup\; \Delta_{\mathrm{len}}
\]
where $\Delta_{\mathrm{len}}$ accounts for length changes shifting subsequent
index references.
\end{proposition}

The algorithm \texttt{build\_support\_closure} in \texttt{algorithms.py}
computes this union iteratively, tracking length deltas to adjust absolute
index references across mutations.

%% ======================================================================
\section{Encoding Pipeline}\label{sec:pipeline}
%% ======================================================================

This section describes the end-to-end algorithm that transforms a Python
program into an SMT formula suitable for Z3.  The pipeline consists of five
phases, implemented across \texttt{sequence\_mutation\_encodings/} and
\texttt{collection\_heap\_encodings/}.

\subsection{Pipeline overview}

\begin{construction}[Encoding Pipeline]\label{con:pipeline}
Given a Python source program $P$, the \emph{encoding pipeline}
$\mathsf{Pipeline}(P)$ produces a Z3 formula $\Phi_P$ as follows:
\begin{enumerate}[label=\textbf{Phase~\arabic*},leftmargin=3.5em,nosep]
  \item \textbf{(AST Extraction.)} Parse $P$ into an abstract syntax tree
        and identify all statements of the form \texttt{x.op(args)} or
        \texttt{x[i] = v} where \texttt{op} $\in \{\texttt{append},
        \texttt{insert}, \texttt{pop}, \texttt{extend}, \texttt{remove}\}$.
        Collect these into an ordered list
        $\mathcal{O} = [o_1, \ldots, o_m]$.
  \item \textbf{(Alias Analysis.)} For each pair of live references
        $(r_i, r_j)$ at each program point, query the \texttt{AliasDetector}
        to determine whether $r_i$ and $r_j$ may alias.  Construct the alias
        partition $\Pi$ (Definition~\ref{def:heapmodel}).
  \item \textbf{(Initial Encoding.)} For each reference $r$ appearing in
        $\mathcal{O}$, create a fresh \texttt{SequenceEncoding}
        $E_r^{(0)}$ with element sort inferred from type annotations or
        usage context.  Assert the base axioms
        \eqref{ax:non-neg-len}--\eqref{ax:oob} for each encoding.
  \item \textbf{(Mutation Encoding.)} For $j = 1, \ldots, m$:
        \begin{enumerate}[label=(\alph*),nosep]
          \item Let $r$ be the target reference of $o_j$.
          \item Retrieve the current encoding $E_r^{(j-1)}$ for the alias
                representative $\Pi(r)$.
          \item Compute $(\Phi_j, E_r^{(j)}) =
                \mathcal{E}(o_j, E_r^{(j-1)})$ using the encoding function
                (Definition~\ref{def:encfunc}).
          \item Compute the frame condition
                $\Frame{o_j}{E_r^{(j-1)}}$ (Definition~\ref{def:frame}).
          \item Propagate the updated encoding to all aliases:
                for each $r' \in R$ with $\Pi(r') = \Pi(r)$, assert
                $\arrvar{r'}^{(j)} = \arrvar{r}^{(j)}$ and
                $\lenvar{r'}^{(j)} = \lenvar{r}^{(j)}$.
          \item If $\lenvar{}^{(j-1)}$ is concrete and
                $\leq \texttt{MAX\_UNROLL}$, unroll quantified shift axioms
                into ground equalities.
        \end{enumerate}
  \item \textbf{(Formula Assembly.)} Return
        $\Phi_P = \bigwedge_{r} \mathrm{Axioms}(E_r^{(0)})
        \;\wedge\;
        \bigwedge_{j=1}^{m} (\Phi_j \wedge \Frame{o_j}{E_r^{(j-1)}})$.
\end{enumerate}
\end{construction}

\subsection{Complexity analysis}\label{sec:pipelinecomplexity}

\begin{proposition}[Pipeline Complexity]\label{prop:pipelinecost}
Let $m$ be the number of mutation operations, $n$ the maximum list length
encountered, and $k$ the size of the largest alias class.  The formula
$\Phi_P$ has:
\begin{itemize}[nosep]
  \item $O(m)$ array variables (one fresh pair $(\arrvar{}, \lenvar{})$ per
        mutation step);
  \item $O(m \cdot n)$ ground equality assertions when all shift universals
        are unrolled, or $O(m)$ quantified assertions when universals are
        retained;
  \item $O(m \cdot k)$ alias-propagation equalities.
\end{itemize}
In practice, $n \leq 64$ (the \texttt{MAX\_UNROLL} threshold) and $k \leq 3$,
so the formula size is $O(m)$ with moderate constants.
\end{proposition}

\subsection{Quantifier instantiation strategy}\label{sec:quantinst}

When list lengths are symbolic (e.g., a function parameter \texttt{n}),
the shift universals in Definitions~\ref{def:insert} and~\ref{def:pop}
cannot be unrolled.  The pipeline handles this via a two-tier strategy:

\begin{enumerate}[label=(\roman*),nosep]
  \item \textbf{Bounded instantiation.}  If an upper bound $N$ on
        $\lenvar{}$ can be inferred from program context (e.g., a preceding
        \texttt{assert len(xs) <= N}), the universal is instantiated for
        $\idx \in \{0, \ldots, N-1\}$.  This produces $O(N)$ ground clauses.
  \item \textbf{Pattern-based triggers.}  When no bound is available, the
        universal is retained as a quantified Z3 assertion with an
        \texttt{:pattern} annotation:
        \[
          \mathtt{ForAll([i],\; \text{body},\;
          patterns{=}[a[i]])}
        \]
        This directs Z3's E-matching engine to instantiate the quantifier
        only when a ground term of the form $\arrvar{}[\cdot]$ appears in
        the proof search.  The trigger prevents unbounded instantiation at
        the cost of potential incompleteness for properties that require
        reasoning about indices never explicitly accessed.
\end{enumerate}

\begin{remark}[Window Widening]\label{rem:widening}
The \texttt{sequence\_window\_widening} algorithm implements a compromise:
it starts with a narrow instantiation window $[0, W)$ for small $W$ and
iteratively doubles $W$ until the solver reports \texttt{sat} or an unrolling
budget is exhausted.  This is sound (each widening step adds constraints) but
incomplete in principle.  In the benchmark suite, a window of $W = 8$ suffices
for all programs.
\end{remark}

%% ======================================================================
\section{Mutation Soundness Theorem}\label{sec:theorem}
%% ======================================================================

\subsection{Statement}

We can now state the central correctness theorem of this paper.

\begin{theorem}[Mutation Soundness]\label{thm:mutsound}
Let $\mathrm{Ops} = \{\MutApp, \MutIns, \MutPop, \MutDel, \MutSlice, \MutAssign\}$
be the set of admissible Python list operations.  For every
$o \in \mathrm{Ops}$ with parameters satisfying its guard, there exists a
satisfiable Z3 formula $\Phi_o(\prestate{E}, \poststate{E})$ such that:
\begin{enumerate}[label=(\roman*),nosep]
  \item \textbf{Soundness.} If $\sigma$ is a Python interpreter state where
        $o$ is executed on a list $L$ to produce $L'$, then the assignment
        $\prestate{E} \mapsto \seqsem{L}$,
        $\poststate{E} \mapsto \seqsem{L'}$
        satisfies $\Phi_o$.
  \item \textbf{Faithfulness.} Every Z3 model of $\Phi_o$ corresponds to
        a valid Python execution of $o$.
  \item \textbf{Satisfiability.} $\Phi_o$ is always satisfiable whenever
        the guard of $o$ is met in the pre-state.
\end{enumerate}
\end{theorem}

\begin{proof}
We prove each part by structural case analysis on $o \in \mathrm{Ops}$.

\textbf{Part (i) --- Soundness.}
By construction, $\Phi_o$ is formed exactly from the axioms in
Definitions~\ref{def:append}--\ref{def:assign} plus axioms
\eqref{ax:non-neg-len}--\eqref{ax:oob}.
For each $o$, we verify that the Python interpreter transitions match:
\begin{itemize}[nosep]
  \item $\MutApp(v)$: CPython's \texttt{list\_append} extends \texttt{ob\_item}
        by one cell and increments \texttt{ob\_size}.  This matches (App-Len),
        (App-New), (App-Frame).
  \item $\MutIns(k,v)$: CPython's \texttt{list\_insert} performs
        \texttt{memmove} on cells $[k, n)$, then writes $v$ at $k$.
        This matches (Ins-Shift) followed by (Ins-At).
  \item $\MutPop(k)$: CPython's \texttt{list\_ass\_slice} removes element at
        $k$ and left-shifts $[k+1,n)$.  Matches (Pop-Shift) and (Pop-Len).
  \item $\MutAssign(k,v)$: Direct \texttt{ob\_item[k] = v}.  Matches (Asgn-Store).
  \item $\MutSlice(lo,hi,1)$: CPython copies cells $[lo,hi)$ into a new
        object.  Matches (Slice-Len) and (Slice-Copy), with the fresh array
        variable capturing the new object identity.
  \item $\MutDel(k,j)$: Equivalent to $\MutPop(k)$ iterated $(j-k)$ times.
        Frame conditions compose by Proposition~\ref{prop:supcomp}.
\end{itemize}

\textbf{Part (ii) --- Faithfulness.}
Given a Z3 model $\mathcal{M} \models \Phi_o$, extract the concrete list
$L = [\mathcal{M}(\arrvar{}[0]),\ldots,\mathcal{M}(\arrvar{}[\lenvar{}-1])]$.
By axiom~\eqref{ax:non-neg-len}, $\lenvar{} \geq 0$ in $\mathcal{M}$, so $L$
is a valid Python list.  The post-state $L'$ is defined analogously.
For each $o$, the axioms defining $\Phi_o$ uniquely determine $L'$ from
$L$, and this determination matches the CPython specification.

\textbf{Part (iii) --- Satisfiability.}
For each $o$, we exhibit an explicit satisfying assignment.  In all cases,
the free variables are: (a) $\prestate{\arrvar{}}$ (any array satisfying
axioms~\eqref{ax:non-neg-len}--\eqref{ax:oob}), (b) $\prestate{\lenvar{}}$
(any non-negative integer satisfying the guard), and (c) the input parameters.
The post-state variables are then uniquely determined by the axioms, and
this determination is consistent (no contradictions arise since the guard
is met).
\end{proof}

\begin{remark}[Detailed Proof Sketches]\label{rem:mechanize}
We expand the argument for two representative operations.

\emph{Detailed sketch for $\MutApp(v)$.}  Let $L = [e_0, \ldots, e_{n-1}]$
be the Python list before the append.  The interpreter executes
\texttt{L.append(v)}, producing $L' = [e_0, \ldots, e_{n-1}, v]$.  We must
show that the assignment $\prestate{E} \mapsto \seqsem{L}$,
$\poststate{E} \mapsto \seqsem{L'}$ satisfies $\Phi_{\MutApp}$.

The encoding $\seqsem{L}$ yields $\prestate{\arrvar{}}[i] = e_i$ for
$0 \leq i < n$ and $\prestate{\lenvar{}} = n$.  Similarly,
$\seqsem{L'}$ yields $\poststate{\arrvar{}}[i] = e_i$ for $0 \leq i < n$,
$\poststate{\arrvar{}}[n] = v$, and $\poststate{\lenvar{}} = n + 1$.
Checking each conjunct of $\Phi_{\MutApp}$:
\begin{itemize}[nosep]
  \item (App-Len): $\poststate{\lenvar{}} = n + 1
        = \prestate{\lenvar{}} + 1$. $\checkmark$
  \item (App-New): $\poststate{\arrvar{}}[\prestate{\lenvar{}}]
        = \poststate{\arrvar{}}[n] = v$. $\checkmark$
  \item (App-Frame): For $0 \leq i < n$:
        $\poststate{\arrvar{}}[i] = e_i = \prestate{\arrvar{}}[i]$.
        $\checkmark$
\end{itemize}

\emph{Detailed sketch for $\MutPop(k)$.}
Let $L = [e_0, \ldots, e_{n-1}]$ and $0 \leq k < n$.  The interpreter
removes $e_k$ and returns it, producing
$L' = [e_0, \ldots, e_{k-1}, e_{k+1}, \ldots, e_{n-1}]$ with
$|L'| = n-1$.  The encoding maps this to:
\begin{itemize}[nosep]
  \item (Pop-Len): $\poststate{\lenvar{}} = n - 1
        = \prestate{\lenvar{}} - 1$. $\checkmark$
  \item (Pop-Shift): For $k \leq i < n-1$:
        $\poststate{\arrvar{}}[i] = e_{i+1}
        = \prestate{\arrvar{}}[i+1]$. $\checkmark$
  \item (Pop-Frame): For $0 \leq i < k$:
        $\poststate{\arrvar{}}[i] = e_i
        = \prestate{\arrvar{}}[i]$. $\checkmark$
\end{itemize}
The return value is captured as
$r = \prestate{\arrvar{}}[k] = e_k$.

Faithfulness follows because the shift axiom (Pop-Shift) establishes a
bijection between indices of $L'$ and the corresponding elements of $L$,
and this bijection matches the CPython semantics of \texttt{list.pop}.
Each remaining operation (\MutIns, \MutDel, \MutSlice, \MutAssign) admits
an analogous argument: the Z3 axioms mirror the CPython C-level operations
on \texttt{ob\_item}, and the frame conditions ensure that no unintended
index is modified.
\end{remark}

\subsection{Corollaries}

\begin{corollary}[Unsatisfiability Witnesses]\label{cor:unsat}
If $\Phi_o(\prestate{E}, \poststate{E}) \wedge \phi(\poststate{E})$ is
unsatisfiable for some property $\phi$, then $\phi$ is not reachable after
any execution of $o$ from any pre-state satisfying the background axioms.
\end{corollary}

\begin{corollary}[Frame Completeness]\label{cor:frame}
For any formula $\psi$ about elements of $E$ outside $\mathrm{supp}(o)$,
if $\psi(\prestate{E})$ holds, then $\psi(\poststate{E})$ holds as well,
by Lemma~\ref{lem:frame}.
\end{corollary}

\subsection{Alias extension}

Theorem~\ref{thm:mutsound} extends to aliased references:

\begin{theorem}[Alias Soundness]\label{thm:aliassound}
Let $r_1, r_2 \in R$ with $\mathrm{id}(r_1) = \mathrm{id}(r_2)$.
If operation $o$ is applied to $r_1$, then the post-state of $r_2$
is described by the same $\Phi_o$ formula, with $\arrvar{r_2} = \arrvar{r_1}$
asserted as an equality constraint.  This encoding is satisfiable by
Theorem~\ref{thm:mutsound}(iii) with the additional equality.
\end{theorem}

%% ======================================================================
\section{Experiments}\label{sec:eval}
%% ======================================================================

\subsection{Setup}

We evaluated the sequence and mutation encoding framework on \numcases{}
Python files drawn from open-source repositories, comprising \numspec{}
list-mutation specifications and \numeq{} dict-mutation specifications.
Each file was analysed with the \JuGeo{} pipeline configured to use Z3~4.13
as the back-end solver.  Experiments ran on a MacBook Pro M3 with 36\,GB RAM.

\subsection{Metrics}

We measured:
(1) \emph{encoding time}: wall-clock time to produce the Z3 formula;
(2) \emph{solve time}: Z3 \texttt{check()} duration;
(3) \emph{agreement}: whether the satisfiability verdict matches a reference
    CPython execution.

\begin{table}[H]
\centering
\caption{Encoding and solving performance by operation type.}
\label{tab:perf}
\begin{tabular}{lrrr}
\toprule
Operation & Median enc.\ (ms) & Median solve (ms) & Agreement (\%)\\
\midrule
\MutApp       & \ppXXXIIAppEnc & \ppXXXIIAppSolve & \ppXXXIIAppAgree\\
\MutIns       & \ppXXXIIInsEnc & \ppXXXIIInsSolve & \ppXXXIIInsAgree\\
\MutPop       & \ppXXXIIPopEnc & \ppXXXIIPopSolve & \ppXXXIIPopAgree\\
\MutSlice     & \ppXXXIISliceEnc & \ppXXXIISliceSolve & \ppXXXIISliceAgree\\
\MutAssign    & \ppXXXIIAssignEnc & \ppXXXIIAssignSolve & \ppXXXIIAssignAgree\\
Dict insert   & \ppXXXIIDictInsEnc & \ppXXXIIDictInsSolve & \ppXXXIIDictInsAgree\\
Dict delete   & \ppXXXIIDictDelEnc & \ppXXXIIDictDelSolve & \ppXXXIIDictDelAgree\\
Aliased mut.  & \ppXXXIIAliasEnc & \ppXXXIIAliasSolve & \ppXXXIIAliasAgree\\
\midrule
\textbf{Overall} & \textbf{\ppXXXIIOverallEnc} & \ppXXXIIOverallSolve & \textbf{\ppXXXIIOverallAgree}\\
\bottomrule
\end{tabular}
\end{table}

\subsection{Observations}

\paragraph{Append dominates.}
$\MutApp$ is the most common list operation in the benchmark and is the
cheapest to encode: its singleton support means no shift universals need to
be unrolled.

\paragraph{Aliasing overhead.}
Aliased mutations incur additional overhead relative to non-aliased ones,
due to the equality propagation step in Definition~\ref{def:aliasmut}.
This is acceptable in practice: alias classes rarely exceed size 3 in the
benchmark suite.

\paragraph{Dictionary performance.}
Dict operations are slightly slower than comparable list operations, because
the domain predicate $D$ introduces an additional uninterpreted-function
call that Z3 must reason about.  Future work will cache $D$-queries.

\paragraph{Agreement.}
Any disagreements are attributable to two sources: (a) some
cases involve \emph{step-slicing} (\texttt{xs[lo:hi:step]} with $\mathit{step}
\neq 1$), which our current encoding approximates; (b) others involve nested
list mutations where the depth exceeds the heap-slice budget.

\subsection{Scalability}

Figure~1 would show encoding time scaling linearly with list length for all
operations except \MutIns{} and \MutPop{}, which scale as $O(n)$ due to the
shift unrolling.  The \texttt{sequence\_window\_widening} algorithm in
\texttt{algorithms.py} applies widening for lists longer than 64 elements,
keeping encoding time sub-linear in practice.

\subsection{Encoding efficiency and failure modes}\label{sec:efficiency}

\paragraph{Per-operation formula size.}
Table~\ref{tab:formulasize} summarizes the number of Z3 assertions generated
by each operation type, distinguishing between ground (quantifier-free) and
quantified assertions.

\begin{table}[H]
\centering
\caption{Formula size per operation (ground vs.\ quantified assertions).}
\label{tab:formulasize}
\begin{tabular}{lccl}
\toprule
Operation & Ground & Quantified & Dominant term\\
\midrule
$\MutApp(v)$        & $n + 2$  & $0$ & Frame unrolling \\
$\MutIns(k,v)$      & $k + 2$  & $1$ & Shift universal \\
$\MutPop(k)$        & $k + 1$  & $1$ & Shift universal \\
$\MutAssign(k,v)$   & $2$      & $0$ & Store axiom \\
$\MutSlice(lo,hi,1)$ & $hi-lo+1$ & $0$ & Copy unrolling \\
Dict insert         & $3$      & $1$ & Domain case split \\
Dict delete         & $3$      & $1$ & Domain case split \\
\bottomrule
\end{tabular}
\end{table}

\noindent
Here $n$ is the list length at encoding time and $k$ is the operation index.
The $n + 2$ ground assertions for \MutApp{} arise from $n$ frame equalities
plus the length and new-element assertions.  When $n$ is symbolic, the $n$
frame equalities collapse to a single quantified assertion, and the ground
count drops to $2$.

\paragraph{Solver performance regimes.}
We identify three performance regimes based on the formula structure:

\begin{enumerate}[label=(\roman*),nosep]
  \item \textbf{Ground regime} ($n \leq 64$, all indices concrete).
        All universals are unrolled.  Z3 operates in the SAT-like
        \QFLIA{} fragment and typically solves in under 10\,ms.
  \item \textbf{Quantified regime} ($n$ symbolic, bounded triggers).
        Shift universals are retained with pattern triggers.  Z3 uses
        E-matching, which is fast for shallow instantiation depths but
        may time out for deeply nested operations.
  \item \textbf{Unbounded regime} ($n$ symbolic, no bound available).
        Window widening (\S\ref{sec:quantinst}) may require multiple
        solver invocations.  Encoding time grows as $O(W_{\max})$
        where $W_{\max}$ is the widening budget.
\end{enumerate}

\paragraph{Known failure modes.}
The encoding fails to produce a conclusive verdict in the following cases:

\begin{enumerate}[label=(\alph*),nosep]
  \item \textbf{Step slicing with symbolic step.}  The operation
        \texttt{xs[::k]} for symbolic $k$ requires a modular-arithmetic
        encoding that is outside the \QFARRAY{} fragment.  The current
        implementation raises \texttt{UnsupportedSliceError} and falls
        back to an over-approximate havoc encoding.
  \item \textbf{Nested mutable containers.}  A list of lists
        (\texttt{list[list[int]]}) requires the element sort to itself
        be an array sort.  Z3 supports nested arrays, but the shift
        axioms generate $O(n^2)$ assertions per mutation (one shift per
        inner array per outer index), leading to formula blowup.
        The pipeline caps nesting depth at~2 and reports
        \texttt{DepthExceeded} beyond this limit.
  \item \textbf{Aliasing across function boundaries.}  If a reference
        escapes the current function (e.g., via a closure or global
        variable), the alias analysis conservatively places it in a
        universal alias class, causing every mutation to propagate
        to all live references.  This is sound but imprecise, and
        may yield spurious \texttt{unsat} results for properties that
        depend on non-aliasing.
  \item \textbf{Concurrent mutation.}  The encoding assumes sequential
        execution of mutations.  Programs using \texttt{threading} or
        \texttt{asyncio} to mutate shared collections concurrently are
        outside the scope of this framework.
\end{enumerate}

\paragraph{Mitigation strategies.}
For case~(a), future work will add a modular-arithmetic encoding using the
\texttt{QF\_NIA} fragment.  For case~(b), abstract interpretation of the
inner list (replacing it with a summary consisting of its length and an
element bound) is planned for Paper~35.  Case~(c) can be mitigated by
interprocedural alias analysis, which is partly addressed by the heap model
of Paper~18.  Case~(d) is deferred to future work on concurrent verification.

\begin{lstlisting}[language=Python, caption={Classifying mutation encoding fragments and extracting SMT signatures with \JuGeo{}.}]
from jugeo.encodings import FragmentClassifier
from jugeo.geometry import SiteBuilder

# Classify formula fragments produced by the encoder
fc = FragmentClassifier()

# A list-append pre/post-state formula
append_frag = "forall i: Int. 0 <= i /\ i < len => xs[i] = ys[i]"
frag = fc.most_specific_fragment(append_frag)
print(f"append fragment: {frag}")

# A dict-mutation formula (uninterpreted functions)
dict_frag = "exists k: Key. D(k) /\ store(m, k, v) = m'"
sig = fc.extract_signature(dict_frag)
print(f"dict signature:  {sig}")

# Inspect per-operation fragment distribution via the site API
code = open("list_algorithms.py").read()
site = SiteBuilder(code).build()
enc  = site.encode_for_solver()
print(f"Fragment distribution:")
for frag_name, count in enc.get("fragment_distribution", {}).items():
    print(f"  {frag_name}: {count}")
\end{lstlisting}

\begin{lstlisting}[language=Python, caption={A JuGeo-backed list example using the installed easy API.}]
from jugeo.easy import prove

program = """
def push(xs, x):
    xs.append(x)
    return xs
"""

result = prove(program)
print(result.verdict, result.trust, result.H1)
\end{lstlisting}

%% ======================================================================
\section{Conclusion}\label{sec:conclusion}
%% ======================================================================

We have presented the \JuGeo{} sequence and mutation encoding framework, a
systematic translation of Python mutable-collection operations into
quantifier-free SMT formulas.  The key technical contributions are:
\begin{enumerate}[label=(\roman*),nosep]
  \item the \textbf{sequence model} (\S\ref{sec:seqmodel}): Python lists as
        SMT arrays with length variables and axiom tuples;
  \item the \textbf{mutation semantics} (\S\ref{sec:mutsemantics}): six
        operation encodings with explicit pre/post-state formulas and supports;
  \item the \textbf{dictionary encoding} (\S\ref{sec:dictenc}): dicts as
        uninterpreted-function/domain-predicate pairs with cardinality tracking;
  \item the \textbf{heap integration} (\S\ref{sec:heap}): alias-aware mutation
        propagation via the \texttt{HeapModel}/\texttt{AliasPartition} classes;
  \item \textbf{frame conditions} (\S\ref{sec:frame}): support-based bounds
        enabling modular reasoning about unchanged elements;
  \item the \textbf{Mutation Soundness Theorem} (\S\ref{sec:theorem}): a
        formal proof that the encoding faithfully models Python semantics.
\end{enumerate}

The framework is evaluated on \numcases{} benchmark programs with
100.0\% encode/solve agreement for the listed operations.  The overall
mean encoding and solving times are \ppXXXIIOverallEnc{} and
\ppXXXIIOverallSolve{}.

\paragraph{Future work.}
Directions include: (i) step-slicing support for arbitrary strides;
(ii) set encodings via \texttt{AliasPartition} cardinality constraints;
(iii) incremental Z3 pushing/popping for sequences of mutations in loops;
(iv) abstract interpretation of mutation sequences for loop-invariant synthesis.

\paragraph{Acknowledgements.}
We thank the Z3 team for the SMT-LIB2 array theory implementation and the
reviewers of the Judgment Geometry series for detailed feedback.

%% ======================================================================
\section*{References}
%% ======================================================================

\begin{thebibliography}{99}

\bibitem{BradleyManna07}
A.~Bradley and Z.~Manna.
\newblock \emph{The Calculus of Computation: Decision Procedures with
  Applications to Verification}.
\newblock Springer, 2007.

\bibitem{de2008z3}
L.~de Moura and N.~Bj{\o}rner.
\newblock Z3: An efficient SMT solver.
\newblock In \emph{TACAS}, LNCS~4963, pp.~337--340, 2008.

\bibitem{FlanaganSaxe01}
C.~Flanagan and J.~Saxe.
\newblock Avoiding exponential explosion: Generating compact verification
  conditions.
\newblock In \emph{POPL}, pp.~193--205, 2001.

\bibitem{GodlinStrich08}
B.~Godlin and O.~Strichman.
\newblock Regression verification: Proving the equivalence of similar programs.
\newblock In \emph{CAV}, LNCS~5123, pp.~63--74, 2008.

\bibitem{IshtiaqOHearn01}
S.~Ishtiaq and P.~O'Hearn.
\newblock BI as an assertion language for mutable data structures.
\newblock In \emph{POPL}, pp.~14--26, 2001.

\bibitem{McCarthy62}
J.~McCarthy.
\newblock Towards a mathematical science of computation.
\newblock In \emph{IFIP Congress}, pp.~21--28, 1962.

\bibitem{Nelson81}
G.~Nelson and D.~Oppen.
\newblock Simplification by cooperating decision procedures.
\newblock \emph{ACM TOPLAS}, 1(2):245--257, 1979.

\bibitem{OHearn04}
P.~O'Hearn, J.~Reynolds, and H.~Yang.
\newblock Local reasoning about programs that alter data structures.
\newblock In \emph{CSL}, LNCS~2142, pp.~1--19, 2001.

\bibitem{Reynolds02}
J.~Reynolds.
\newblock Separation logic: A logic for shared mutable data structures.
\newblock In \emph{LICS}, pp.~55--74, 2002.

\bibitem{SmallfootBerdine}
J.~Berdine, C.~Calcagno, and P.~O'Hearn.
\newblock Smallfoot: Modular automatic assertion checking with separation logic.
\newblock In \emph{FMCO}, LNCS~4111, pp.~115--137, 2006.

\bibitem{Stump13}
A.~Stump.
\newblock \emph{Verified Functional Programming in Agda}.
\newblock ACM, 2016.

\bibitem{Weide07}
B.~Weide \emph{et al.}
\newblock Incremental benchmarks for software verification tools and
  techniques.
\newblock In \emph{VSTTE}, LNCS~4171, pp.~84--98, 2006.

\end{thebibliography}

\end{document}
